\subsection{伴随表示}

设 \(G\) 是一个李群。考虑解析左运算律

\[
(g, g ^ {\prime}) \mapsto g g ^ {\prime} g ^ {- 1} = (\operatorname{Int} g) g ^ {\prime}
\]

从 G 到 G。根据第3号，此运算律定义了一个从 \(\mathcal{T}^{(\infty)}(\mathrm{G})\times \mathcal{T}^{(\infty)}(\mathrm{G})\) 到 \(\mathcal{T}^{(\infty)}(\mathrm{G})\) 的双线性映射，本号中记为 \(\top\)。根据第3号命题13，

\[
(t * t ^ {\prime}) \top t ^ {\prime \prime} = t \top (t ^ {\prime} \top t ^ {\prime \prime})\tag{22}
\]

对任意 \(t, t', t''\)，它们属于 \(\mathcal{T}^{(\infty)}(\mathrm{G})\)。根据第3号命题14 (i)，

\[
\varepsilon_ {g} \top t = (\operatorname{Int} g) _ {*} t\tag{23}
\]

对任意 \(g \in G\) 及 \(t \in \mathcal{T}^{(\infty)}(G)\)。特别地，映射 \(t \mapsto \varepsilon_{g} \top t\) 是从 \(\mathcal{T}^{(\infty)}(G)\) 到 \(\mathcal{T}^{(\infty)}(G)\) 的双代数 \(\mathcal{T}^{(\infty)}(G)\) 的自同构。它在 \(U(G)\)、\(U_{s}(G)\)、\(L(G)\) 上的限制分别记为 \(\mathrm{Ad}_{\mathrm{U(G)}}(g)\)、\(\mathrm{Ad}_{\mathrm{U_s(G)}}(g)\)、\(\mathrm{Ad}_{\mathrm{L(G)}}(g)\)。不致混淆时，我们常用 \(\mathrm{Ad}(g)\) 代替 \(\mathrm{Ad}_{\mathrm{L(G)}}(g)\)。由 (23)，\(\mathrm{Ad}(g)\) 是在 \(e\) 处的 \(\operatorname{Int}(g)\) 的切映射。它是可赋范

李代数 \(\mathbf{L}(\mathbf{G})\) 的自同构。当 \(\mathbf{K}\) 的特征为 0 时，\(\operatorname{Ad}_{\mathbf{U}(\mathbf{G})}(g)\) 是 \(\mathbf{U}(\mathbf{G})\) 的唯一自同构，它延拓 \(\operatorname{Ad}(g)\)。

若 \(\phi\) 是从李群 G 到李群 H 的态射，则

\[
\phi_ {*} (t \top t ^ {\prime}) = \phi_ {*} (t) \top \phi_ {*} (t ^ {\prime})\tag{24}
\]

对任意 \(t, t'\)，它们属于 \(\mathcal{T}^{(\infty)}(\mathrm{G})\)；这由第3号命题15推出。

命题 43. 设 \(t, u\) 属于 \(\mathcal{T}^{(\infty)}(\mathrm{G})\)。设 \(\sum_{i=1}^{n} t_i \otimes t'_i\) 是 \(t\) 在余乘法下的像。则

\[
t \top u = \sum_ {i = 1} ^ {n} t _ {i} * u * t _ {i} ^ {\prime \vee}.
\]

按定义，\(t \top u\) 是 \(t \otimes u\) 在映射 \((g, g') \mapsto gg'g^{-1}\) 下的像，该映射从 \(G \times G\) 到 \(G\)。现在，这个映射由下列映射复合得到：

\[
\begin{array}{l l} \alpha \colon (g, g ^ {\prime}) \mapsto (g, g, g ^ {\prime}) & \text {of G \times G into G \times G \times G} \\ \beta \colon (g, g ^ {\prime}, g ^ {\prime \prime}) \mapsto (g, g ^ {\prime - 1}, g ^ {\prime \prime}) & \text {of G \times G \times G into G \times G \times G} \\ \gamma \colon (g, g ^ {\prime}, g ^ {\prime \prime}) \mapsto g g ^ {\prime \prime} g ^ {\prime} & \text {of G \times G \times G into G}. \end{array}
\]

另一方面：

\[
\begin{array}{c} \alpha_ {*} (t \otimes u) = \sum_ {i = 1} ^ {n} (t _ {i} \otimes t _ {i} ^ {\prime}) \otimes u = \sum_ {i = 1} ^ {n} t _ {i} \otimes t _ {i} ^ {\prime} \otimes u \\ \beta_ {*} (\sum_ {i = 1} ^ {n} t _ {i} \otimes t _ {i} ^ {\prime} \otimes u) = \sum_ {i = 1} ^ {n} t _ {i} \otimes t _ {i} ^ {\prime \vee} \otimes u \\ \gamma_ {*} (\sum_ {i = 1} ^ {n} t _ {i} \otimes t _ {i} ^ {\prime \vee} \otimes u) = \sum_ {i = 1} ^ {n} t _ {i} * u * t _ {i} ^ {\prime \vee}. \end{array}
\]

推论 1. 设 \(u \in \mathrm{L}(\mathbf{G})\)，且 \(u' \in \mathcal{T}^{(\infty)}(\mathbf{G})\)。则 \(u \top u' = u * u' - u' * u\)。

u 在余乘法下的像是 \(u \otimes \varepsilon_{e} + \varepsilon_{e} \otimes u\)，所以

\[
u \top u ^ {\prime} = u * u ^ {\prime} * \varepsilon_ {e} + \varepsilon_ {e} * u ^ {\prime} * u ^ {\vee} = u * u ^ {\prime} - u ^ {\prime} * u.
\]

推论 2. 设 \(t \in \mathcal{T}^{(\infty)}(\mathrm{G})\) 且 \(g \in \mathrm{G}\)。则 \(\varepsilon_g \top t = \varepsilon_g * t * \varepsilon_g^{-1}\)。若 \(t \in \mathrm{L}(\mathrm{G})\)，则 \(\varepsilon_g \top t = gtg^{-1}\)（后一个乘积在群 \(\mathrm{T}(\mathrm{G})\) 中计算）。

\(\varepsilon_{g}\) 在余乘法下的像是 \(\varepsilon_{g} \otimes \varepsilon_{g}\)。

推论 3. 设 \(a \in \mathrm{L}(\mathrm{G})\)。由 \(a\) 及左运算 \(g \mapsto \operatorname{Int} g\) 定义的 \(\mathrm{G}\) 在 \(\mathrm{G}\) 上的向量场是场 \(\mathbf{R}_a - \mathbf{L}_a\)。

这个场在 \(g\) 处的值为

\[
\begin{array}{r l} a \top \varepsilon_ {g} = a * \varepsilon_ {g} - \varepsilon_ {g} * a & (\text { Corollary   1 }) \\ = (R _ {a}) _ {g} - (L _ {a}) _ {g} & (\text { Definition   5 }). \end{array}
\]

对任意 \(g \in G\) 及 \(t \in \mathrm{L}(\mathrm{G})\)，

\[
(\mathrm{Ad} g) (t) = \varepsilon_ {g} \top t = \varepsilon_ {g} * t * \varepsilon_ {g ^ {- 1}} = g t g ^ {- 1}.\tag{25}
\]

由于 Ad \(g = \mathrm{T}_{e}(\mathrm{Int} g)\)，第11号命题42证明了 Ad 是 G 在可赋范空间 \(\mathrm{L}(\mathrm{G})\) 上的解析线性表示。

定义 8. G 在 L(G) 上的表示 Ad 称为 G 的伴随表示。

命题 44. 对任意 \(a \in \mathbf{L}(\mathbf{G})\)，

\[
(\mathrm{L} (\mathrm{Ad})) (a) = \mathrm{ad} _ {\mathrm{L} (\mathrm{G})} a.
\]

设 \(b \in \mathrm{L}(\mathrm{G})\)。根据第11号命题42 (ii)及命题43的推论3，

\[
(\mathrm{L} (\mathrm{Ad})) (a). b = - [ \mathrm{R} _ {a} - \mathrm{L} _ {a}, \mathrm{L} _ {b} ] (e).
\]

现在 \(R_{a} \circ L_{b} = L_{b} \circ R_{a}\)（第6号命题23 (ii)），因而 \([R_{a}, L_{b}] = 0\)；再利用命题23 (ii)，

\[
(\mathrm{L} (\mathrm{Ad})) (a). b = [ \mathrm{L} _ {a}, \mathrm{L} _ {b} ] (e) = \mathrm{L} _ {[ a, b ]} (e) = [ a, b ] = (\mathrm{ad} _ {\mathrm{L} (\mathrm{G})} a) b.
\]

命题 45. 假设 G 是有限维的，且 K 的特征为 0。设 s 是满足 \(\geqslant0\) 的整数。则映射 \(\pi: g \mapsto \operatorname{Ad}_{\mathrm{U}_{s}(\mathbf{G})}(g)\) 是 G 在 \(\mathrm{U}_{s}(\mathrm{G})\) 上的解析线性表示，并且 \(\mathrm{L}(\pi)a = \operatorname{ad}_{\mathrm{U}_{s}(\mathbf{G})}a\) 对任意 \(a \in \mathrm{L}(\mathrm{G})\) 成立。

线性表示 \(\pi\) 是 \(\bigoplus_{r=0}^{\infty} \mathrm{T}^{r}(\mathrm{Ad})\) 的一个商，因而是解析的。对于 \(a \in \mathrm{L}(\mathrm{G})\) 以及 \(x_{1}, x_{2}, \ldots, x_{s}\)，它们属于 \(\mathrm{L}(\mathrm{G})\)，

\[
\begin{array}{l l} (\mathrm{L} (\pi) a) (x _ {1} x _ {2} \dots x _ {s}) = \sum_ {i = 1} ^ {s} x _ {1} \dots (\mathrm{L} (\mathrm{Ad}) a. x _ {i}) \dots x _ {s} & \text {(Proposition 41)} \\ = \sum_ {i = 1} ^ {s} x _ {1} \dots ([ a, x _ {i} ]) \dots x _ {s} & \text {(Proposition 44)} \\ = (\mathrm{Ad} _ {\mathrm{U} _ {s} (\mathrm{G})} a) (x _ {1} x _ {2} \dots x _ {s}). \end{array}
\]

命题 46. 设 \(h \in \mathbf{G}\)、\(x \in \mathrm{T}_h(\mathbf{G})\) 且 \(a \in \mathrm{L}(\mathbf{G})\)。设 \(\phi\) 是映射 \((g, g') \mapsto gg'g^{-1}\)，从 \(\mathbf{G} \times \mathbf{G}\) 到 \(\mathbf{G}\)。像 \(y\) 是 \((a, x) \in \mathrm{T}_e(\mathbf{G}) \times \mathrm{T}_h(\mathbf{G})\) 在 \(\mathrm{T}_{(e,h)}(\phi)\) 下的像，且 \(y = x + h((\mathrm{Ad}h^{-1})a - a)\)。

有

\[
\begin{array}{r l} & y = (\mathrm{T} _ {(e, h)} \phi) (a \otimes \varepsilon_ {h} + \varepsilon_ {e} \otimes x) \\ & \quad = a \top \varepsilon_ {h} + \varepsilon_ {e} \top x \\ & \quad = a * \varepsilon_ {h} - \varepsilon_ {h} * a + x \\ & \quad = h ((\mathrm{Ad} h ^ {- 1}) a) - h a + x. \end{array}
\]

命题 47. 设 G 是李群，H 和 E 是 G 的李子群，并假设 \(hEh^{-1} = E\) 对所有 \(h \in H\) 成立。则 \(\mathcal{T}^{(\infty)}(H) \top \mathcal{T}^{(\infty)}(E) \subset \mathcal{T}^{(\infty)}(E)\)。特别地，\(\operatorname{Ad}(H)(L(E)) \subset L(E)\) 且 \([L(H), L(E)] \subset L(E)\)。

若 \(t \in \mathcal{T}^{(\infty)}(\mathrm{H})\) 且 \(t' \in \mathcal{T}^{(\infty)}(\mathrm{E})\)，则 \(t \otimes t' \in \mathcal{T}^{(\infty)}(\mathrm{H} \times \mathrm{E})\)，并且 \(\mathrm{H} \times \mathrm{E}\) 在映射 \((g, g') \mapsto gg'g^{-1}\) 下的像包含于 \(\mathrm{E}\)。

命题 48. 设 G 是李群，H 和 E 是 G 的李子群。假设作为李群，G 是 H 与 E 的半直积。设 \(\rho\) 是李群 G 的线性表示 \(g \mapsto (\mathrm{Ad} g) \mid \mathrm{L}(E)\)，作用在 \(\mathrm{L}(E)\) 上（参见命题47），并设 \(\sigma\) 是 \(\rho\) 在 H 上的限制。则：

\begin{enumerate}
\def\labelenumi{(\roman{enumi})}
\item
  \(\mathbf{L}(\mathbf{G})\) 是 \(\mathbf{L}(\mathbf{H})\) 与 \(\mathbf{L}(\mathbf{E})\) 的拓扑直和；
\item
  L(H) 是 L(G) 的子代数，而 L(E) 是 L(G) 的理想；
\item
  \(\mathbf{L}(\sigma)\) 是 \(\mathbf{L}(\mathbf{H})\) 在 \(\mathbf{L}(\mathbf{E})\) 的导子李代数上的线性表示；
\item
  L(G) 是由 L(σ) 定义的 L(H) 与 L(E) 的半直积（第 I 章 § 1，第8号）。
\item
  显然，且 (ii) 由命题47推出。\(\mathbf{L}(\sigma) = \mathbf{L}(\rho) \mid \mathbf{L}(\mathbf{H})\)。现在根据命题40（第11号）和命题44（第12号），\(\mathbf{L}(\rho)(t)\) 对所有 \(t \in \mathbf{L}(\mathbf{G})\) 都是 \(\mathrm{ad}_{L(\mathbf{G})} t\) 在 \(\mathbf{L}(\mathbf{E})\) 上的限制。这证明了 (iii)。利用 (i) 和 (ii)，也可证明 (iv)。
\end{enumerate}

推论。设 G 是李群。赋予 \(\mathrm{T}_{e}(\mathrm{G})\) 其唯一的交换李代数结构。设 \(\tau\) 是 \(\mathrm{L}(\mathrm{G})\) 的一个伴随表示。则 \(\mathrm{T}(\mathrm{G})\) 的李代数是 \(\mathrm{L}(\mathrm{G})\) 与 \(\mathrm{T}_{e}(\mathrm{G})\) 的半直积，由 \(\tau\) 定义。换言之，对于 \(x, x'\) 属于 \(\mathrm{L}(\mathrm{G})\) 且 \(y, y'\) 属于 \(\mathrm{T}_{e}(\mathrm{G})\)，

\[
[ (x, y), (x ^ {\prime}, y ^ {\prime}) ] = ([ x, x ^ {\prime} ], [ x, y ^ {\prime} ] + [ y, x ^ {\prime} ])
\]

（左侧的括号在 \(\mathbf{L}(\mathbf{T}(\mathbf{G}))\) 中计算，右侧的括号在 \(\mathbf{L}(\mathbf{G})\) 中计算）。

这由命题48及 § 2 第2号的命题6推出。

命题 49. 设 A 是完备可赋范含幺结合代数。将 A 与 L(A*) 识别。则若 \(g \in \mathbf{A}^*\) 且 \(y \in \mathbf{A}\)，(Ad g) \(y = gyg^{-1}\)。

回忆 \(\operatorname{Ad} g = \mathrm{T}_{1}(\operatorname{Int} g)\)。设 \(u_{g}\) 是从 A 到 A 的映射 \(x \mapsto gxg^{-1}\)。\(\mathrm{A}^{*}\) 到 A 的恒等图将 \(\operatorname{Int} g\) 变为 \(u_{g} \mid \mathrm{A}^{*}\)。这个映射在 \(\mathrm{A}^{*}\) 每一点处的切映射都等于 \(u_{g}\)，故得命题。

推论。对任意 \(g \in \mathbf{A}^*\)，设 \(i(g)\) 是由 \(y \mapsto \mathrm{gyg}^{-1}\) 定义的 \(\mathbf{A}\) 的自同构，从而 \(i\) 是 \(\mathbf{A}^*\) 在 \(\mathbf{A}\) 上的解析线性表示。对任意 \(z \in \mathrm{L}(\mathbf{A}^*) = \mathbf{A}\)，\(\mathrm{L}(i)z\) 是由 \(y \mapsto zy - yz\) 定义的 \(\mathbf{A}\) 的内导子。

这由命题49和命题44推出。

